\section{Simulation Parameters and Reproducibility}
\label{app:reproducibility}

This appendix documents the exact parameters, transition matrices, and
algorithms executed by \texttt{scripts/run\_simulations.py} and
\texttt{scripts/generate\_figures.py}.

\subsection{Markov Attention Transition Matrices (Seed \texttt{202601})}
\label{app:rep-matrices}

For $n = 4$ partners $\{G_1, G_2, G_3, G_4\}$, the regime transition matrices
used in \cref{sec:dynamics,sec:simulations} are:
\begin{align*}
  P^{(\mathrm{wd})} &=
  \begin{pmatrix}
    0.62 & 0.18 & 0.14 & 0.06 \\
    0.24 & 0.54 & 0.14 & 0.08 \\
    0.28 & 0.18 & 0.44 & 0.10 \\
    0.30 & 0.20 & 0.15 & 0.35
  \end{pmatrix}, &
  P^{(\mathrm{we})} &=
  \begin{pmatrix}
    0.44 & 0.12 & 0.14 & 0.30 \\
    0.20 & 0.32 & 0.12 & 0.36 \\
    0.22 & 0.14 & 0.28 & 0.36 \\
    0.16 & 0.10 & 0.12 & 0.62
  \end{pmatrix}, \\[6pt]
  P^{(\mathrm{st})} &=
  \begin{pmatrix}
    0.30 & 0.52 & 0.12 & 0.06 \\
    0.14 & 0.72 & 0.08 & 0.06 \\
    0.20 & 0.50 & 0.22 & 0.08 \\
    0.22 & 0.48 & 0.12 & 0.18
  \end{pmatrix}.
\end{align*}
Initial weights are $\mathbf{w}(0) = (0.34, 0.26, 0.20, 0.20)$ starting in slot
$S_0 = 1$, with daily stress probability $p_{\mathrm{st}} = 0.15$ and EMA smoothing
$\alpha = 0.16$.

\subsection{Simulation Algorithms}
\label{app:rep-algos}

\begin{lstlisting}[caption={Pseudocode for social-graph cascade percolation trials (Seed \texttt{202603}).},label={lst:cascade-algo}]
procedure CASCADE_TRIAL(n, q, topology, p_tx=0.72, theta=0.60):
    A <- generate_adjacency(n, q, topology)
    seed <- uniform_int(0, n - 1)
    exposed <- {seed}; queue <- [seed]
    while queue not empty:
        u <- pop_front(queue)
        for v in neighbors(A, u):
            if v not in exposed and uniform_01() < p_tx:
                add(exposed, v)
                push_back(queue, v)
    frac <- |exposed| / n
    return (frac >= theta), frac
\end{lstlisting}

\begin{lstlisting}[caption={Pseudocode for Hermitian unitary-rotation uncertainty sweep (Seed \texttt{202604}).},label={lst:uncertainty-algo}]
procedure UNCERTAINTY_SAMPLE(d, gamma):
    D_R <- diag(1, 2, ..., d)
    D_L <- diag(linspace(-1, 1, d))
    H <- normalized_random_hermitian(d)
    U <- exp(i * gamma * H)
    L <- U^\dagger * D_L * U
    psi <- random_haar_unit_vector(d)
    r <- real(<psi, D_R * psi>); l <- real(<psi, L * psi>)
    Delta_R <- norm((D_R - r * I) * psi)
    Delta_L <- norm((L - l * I) * psi)
    robertson_bound <- 0.5 * abs(imag(<psi, (D_R*L - L*D_R) * psi>))
    cov_term <- real(<(D_R - r*I)*psi, (L - l*I)*psi>)
    schrodinger_bound <- sqrt(cov_term^2 + robertson_bound^2)
    return Delta_R * Delta_L, robertson_bound, cov_term, schrodinger_bound
\end{lstlisting}

\subsection{Full Build and Verification Commands}
\label{app:rep-commands}

From the repository root, the entire pipeline regenerates deterministically via:
\begin{lstlisting}
make figures   # writes data CSVs, summary.json, and PDF/SVG figures
make lean      # runs lake build in lean/ (44 jobs, 20 modules, zero sorry)
make test      # runs python3 -m unittest discover -s tests -v
make pdf       # compiles paper/main.tex to paper/output/
\end{lstlisting}
Running \texttt{make all} executes all four stages sequentially.
